\section*{NS97 audit note: a trace normalization does not change free-step rays}
\paragraph{Classification.}
This is elementary homogeneity applied to NS96, recorded to screen a
suggested continuation. It is not a new arithmetic estimate or a new
research-map theorem. The complete NS96 proof and its prior report are
preserved unchanged.

Use NS96's linear map $T_N$ and old-space projection $P=I-\Pi_N$.
For $\Gamma\succeq0$, set $v=P T_N(\Gamma)$. For a fixed actual residual
$r_N$ and $v\ne0$, the best real line step has gain
\[
 \sup_{t\in\mathbb R}\{2t\langle r_N,v\rangle-t^2\|v\|^2\}
 =\frac{\langle r_N,v\rangle^2}{\|v\|^2}.
\]
For any $a>0$ this quotient is unchanged when $v$ is replaced by $av$.
If $\Gamma\ne0$, PSD implies $g=\operatorname{tr}\Gamma>0$.
Replacing $\Gamma$ by $\Gamma/g$ makes its trace one and replaces $v$
by $v/g$; multiplying the line step by $g$ recovers the same correction.
Consequently a trace-one constraint (or a positive trace upper bound)
does not restrict achievable correction rays when a free step scale is
retained. Even a nonnegative free scale suffices for the unrestricted
PSD-mixture class: NS96 proves its projected cone is a linear space.

For a concrete pair $[i,j]>N$, the two matrices
\[
 \frac12(e_i\pm e_j)(e_i\pm e_j)^T
\]
are PSD of trace one, and project to $\pm\phi_{[i,j]}$.
Thus normalization also fails to force a numerator sign, including for
these single rank-one choices. This does not identify the full single-
rank-one class with the full signed span.

If the trace budget instead constrains the \emph{actual total update}
after the step size is absorbed, the scaling argument does not remove
it. The same is true of a prescribed step size, a genuine restriction on
coefficient ratios, or an independently prescribed residual-dependent
choice. Those are different hypotheses. None is supplied with a cofinal
arithmetic numerator/cost estimate here. A trace bound alone should not
be advertised as a new direction under free line optimization.
